% EDC Book IV: Nuclear Structure from Topological Pinning
% Preamble file — packages and macros
% Date: 2026-02-09

% ============================================================
% PACKAGES
% ============================================================

% Encoding and fonts
\usepackage[utf8]{inputenc}
\usepackage[T1]{fontenc}
\usepackage{lmodern}

% Mathematics
\usepackage{amsmath,amssymb,amsthm}
\usepackage{mathtools}
\usepackage{bm}  % Bold math

% Layout
\usepackage{geometry}
\geometry{
    a4paper,
    left=3cm,
    right=3cm,
    top=3cm,
    bottom=3cm
}

% Headers and footers
\usepackage{fancyhdr}
\pagestyle{fancy}
\fancyhf{}
\fancyhead[LE,RO]{\thepage}
\fancyhead[RE]{\leftmark}
\fancyhead[LO]{\rightmark}
\renewcommand{\headrulewidth}{0.4pt}

% Tables
\usepackage{booktabs}
\usepackage{longtable}
\usepackage{array}
\usepackage{multirow}

% Graphics
\usepackage{graphicx}
\graphicspath{{figures/}}

% Colors and boxes
\usepackage{xcolor}
\usepackage{tcolorbox}
\tcbuselibrary{breakable,skins,theorems}

% Code listings
\usepackage{listings}
\lstset{
    language=Python,
    basicstyle=\ttfamily\small,
    keywordstyle=\color{blue!70!black}\bfseries,
    commentstyle=\color{green!50!black}\itshape,
    stringstyle=\color{red!70!black},
    numbers=left,
    numberstyle=\tiny\color{gray},
    numbersep=5pt,
    breaklines=true,
    breakatwhitespace=true,
    frame=single,
    framesep=3pt,
    xleftmargin=15pt,
    xrightmargin=5pt,
    showstringspaces=false,
    tabsize=4,
    captionpos=b
}

% Cross-references
\usepackage{hyperref}

% Build fingerprint (set by build script, defaults for local builds)
\newcommand{\BuildTag}{local}
\newcommand{\BuildSHA}{unknown}
\newcommand{\BuildDate}{\today}
\newcommand{\BuildPages}{?}
% Load build_info.tex if it exists (generated by tools/build_release_pdf.sh)
\IfFileExists{build_info.tex}{% Auto-generated build info - DO NOT EDIT
\renewcommand{\BuildTag}{book4-narrative-hardened-v1}
\renewcommand{\BuildSHA}{6237e56}
\renewcommand{\BuildDate}{20260212}
}{}

\hypersetup{
    colorlinks=true,
    linkcolor=blue!70!black,
    citecolor=green!50!black,
    urlcolor=blue!80!black,
    bookmarksnumbered=true,
    bookmarksopen=true,
    pdftitle={EDC Book IV: Cluster Structure from Topological Pinning},
    pdfauthor={EDC Research},
    pdfsubject={Build: \BuildTag{} | \BuildSHA{} | \BuildDate{}},
    pdfkeywords={EDC, topological pinning, brane tension, Book IV, \BuildSHA{}}
}
\usepackage{cleveref}

% Bibliography (if needed)
% \usepackage[backend=biber,style=numeric-comp]{biblatex}

% ============================================================
% THEOREM ENVIRONMENTS
% ============================================================

\theoremstyle{plain}
\newtheorem{theorem}{Theorem}[chapter]
\newtheorem{proposition}[theorem]{Proposition}
\newtheorem{lemma}[theorem]{Lemma}
\newtheorem{corollary}[theorem]{Corollary}

\theoremstyle{definition}
\newtheorem{definition}[theorem]{Definition}
\newtheorem{axiom}{Axiom}[chapter]

\theoremstyle{remark}
\newtheorem{remark}[theorem]{Remark}
\newtheorem{note}[theorem]{Note}

% ============================================================
% EDC-SPECIFIC MACROS
% ============================================================

% Brane parameters
\newcommand{\bsigma}{\sigma}           % Brane tension
\newcommand{\bdelta}{\delta}           % 5D length scale
\newcommand{\Lzero}{L_0}               % Characteristic length
\newcommand{\Tstar}{T_*}               % Temperature scale
\newcommand{\Mfive}{M_5}               % 5D Planck mass

% Derived quantities
\newcommand{\Kpin}{K}                  % Pinning constant
\newcommand{\taun}{\tau_n}             % Neutron lifetime
\newcommand{\VB}{V_B}                  % Barrier height
\newcommand{\Earm}{E_{\text{arm}}}     % Per-arm deformation energy (Layer A)
\newcommand{\SE}{S_E}                  % Euclidean action
\newcommand{\omegaz}{\omega_0}         % Oscillation frequency

% Symmetry groups
\newcommand{\Zsix}{\mathbb{Z}_6}
\newcommand{\Zthree}{\mathbb{Z}_3}
\newcommand{\Ztwo}{\mathbb{Z}_2}
\newcommand{\pione}{\pi_1}             % Fundamental group

% Topological
\newcommand{\Sfive}{S^5}               % 5-sphere
\newcommand{\Sone}{S^1}                % Circle

% Coordination
\newcommand{\ncoord}{n}                % Coordination number
\newcommand{\dfrustration}{d(n)}       % Frustration distance

% Units
\newcommand{\MeV}{\,\mathrm{MeV}}
\newcommand{\fm}{\,\mathrm{fm}}
\newcommand{\second}{\,\mathrm{s}}

% ============================================================
% EPISTEMIC TAGS
% ============================================================

% Usage: \epistemictag{Der} or \epistemictag{P}
\newcommand{\epistemictag}[1]{%
    \textsuperscript{\textcolor{blue!70!black}{[#1]}}%
}

\newcommand{\tagDer}{\epistemictag{Der}}    % Derived
\newcommand{\tagDc}{\epistemictag{Dc}}      % Derived with constraints
\newcommand{\tagP}{\epistemictag{P}}        % Postulated
\newcommand{\tagI}{\epistemictag{I}}        % Input
\newcommand{\tagBL}{\epistemictag{BL}}      % Baseline
\newcommand{\tagOpen}{\epistemictag{OPEN}}  % Open problem
\newcommand{\tagOPEN}{\epistemictag{OPEN}}  % Alias (used in several chapters)
\newcommand{\tagCal}{\epistemictag{Cal}}    % Calibrated
\newcommand{\tagM}{\epistemictag{M}}        % Model/Mechanism
\newcommand{\tagDef}{\epistemictag{Def}}    % Definition
\newcommand{\tagDcI}{\epistemictag{Dc+I}}   % Derived with constraints + input

% ============================================================
% BOXES FOR SPECIAL CONTENT
% ============================================================

% Derivation box
\newtcolorbox{derivationbox}[1][]{%
    enhanced,
    colback=blue!5,
    colframe=blue!70!black,
    fonttitle=\bfseries,
    title=Derivation,
    #1
}

% Result box
\newtcolorbox{resultbox}[1][]{%
    enhanced,
    colback=green!5,
    colframe=green!50!black,
    fonttitle=\bfseries,
    title=Result,
    #1
}

% Warning/quarantine box
\newtcolorbox{quarantinebox}[1][]{%
    enhanced,
    colback=red!5,
    colframe=red!70!black,
    fonttitle=\bfseries,
    title=Quarantine: External Comparison,
    #1
}

% Open problem box (renamed to avoid amsthm conflict)
\newtcolorbox{edcopenbox}[1][]{%
    enhanced,
    colback=orange!5,
    colframe=orange!70!black,
    fonttitle=\bfseries,
    title=Open Problem,
    #1
}
% Assumption box
\newtcolorbox{assumptionbox}[1][]{%
    enhanced,
    colback=yellow!5,
    colframe=yellow!50!black,
    fonttitle=\bfseries,
    title=Assumptions,
    #1
}

% Chapter introduction box
\newtcolorbox{chapterbox}[1][]{%
    enhanced,
    colback=gray!5,
    colframe=gray!50!black,
    fonttitle=\bfseries,
    #1
}

% Observer-side projection box
% NOTE: This box contains projection labels (e.g., "proton", "neutron") that map
% 5D EDC objects to 3D/4D observer measurement labels. These are LABELS ONLY,
% not mechanism explanations. The contamination scan ignores observerbox content.
\newtcolorbox{observerbox}[1][]{%
    enhanced,
    breakable,
    colback=purple!5,
    colframe=purple!50!black,
    fonttitle=\bfseries,
    title=Observer-side projection,
    #1
}

% Chapter abstract environment (book class doesn't have this)
\newenvironment{abstract}{%
    \begin{quote}\small\itshape
}{%
    \end{quote}\medskip
}

% ============================================================
% PROVENANCE TRACKING
% ============================================================

% Usage: \source{filename.md}
% NOTE: Source tracking is for editorial use only — output is suppressed in PDF
\newcommand{\source}[1]{}  % Suppressed: internal provenance only

% Usage: \provenance{Ch.6}{INSTANTON_DERIVATION_CHAIN.md}
% NOTE: Suppressed in final PDF
\newcommand{\provenance}[2]{}  % Suppressed: internal provenance only

% ============================================================
% LAYER MARKERS
% ============================================================

\newcommand{\LayerA}{\textcolor{green!50!black}{\textbf{[Layer A]}}}
\newcommand{\LayerB}{\textcolor{orange!70!black}{\textbf{[Layer B]}}}

% ============================================================
% EDC NATIVE ONTOLOGY MACROS
% ============================================================
% Canonical terminology for all EDC books. See ontology/EDC_ONTOLOGY_CANON.md
% RULE: Use these macros in Layer A; do NOT type conventional particle names.

% --- Core Object Macros (Layer A) ---

% Junction states (baryonic)
\newcommand{\AnchorJunction}{anchor junction}
\newcommand{\MetastableJunction}{metastable junction}
\newcommand{\Junction}{junction}
\newcommand{\Constituent}{constituent}

% Loop states (leptonic)
\newcommand{\LoopState}{loop state}
\newcommand{\AntiLoop}{anti-loop}
\newcommand{\LoopExcitation}{loop excitation}

% Mode excitations (bosonic)
\newcommand{\EdgeMode}{edge-mode excitation}
\newcommand{\BulkMode}{bulk mode}

% Cluster states (nuclear)
\newcommand{\ClusterState}{cluster state}
\newcommand{\ClosedFour}{closed-4 unit}
\newcommand{\PinningCluster}{pinning cluster}
\newcommand{\HighCoordCluster}{high-coordination cluster}

% --- Process Macros ---
\newcommand{\JunctionTransition}{junction transition}
\newcommand{\ClosedFourRelease}{closed-4 release}
\newcommand{\Pinning}{pinning}
\newcommand{\Frustration}{frustration}
\newcommand{\BarrierTunneling}{barrier tunneling}

% --- Adjective/Property Macros ---
\newcommand{\HighCoord}{high-coordination}
\newcommand{\ForbiddenZone}{forbidden zone}
\newcommand{\StableCoord}{stable coordination}

% --- Additional Symmetry Macros ---
\newcommand{\Msix}{M_6}                % Coordination lattice

% --- Policy Lock Sentence ---
\newcommand{\EDCPolicyLockSentence}{%
    \textit{This chapter uses EDC-native terminology exclusively. Conventional
    particle names appear only in Appendix~\ref{app:analogies} (analogies) and
    Appendix~\ref{app:quarantine} (calibration). See the Glossary for term
    definitions.}%
}

% --- Analogy-Only Macros (APPENDIX X ONLY!) ---
% WARNING: These macros are for Appendix X translation tables ONLY.
% Do NOT use in Layer A chapters. Grep scan will flag direct use.
\newcommand{\AnalogProton}{proton}            % APPENDIX X ONLY
\newcommand{\AnalogNeutron}{neutron}          % APPENDIX X ONLY
\newcommand{\AnalogElectron}{electron}        % APPENDIX X ONLY
\newcommand{\AnalogNucleus}{nucleus}          % APPENDIX X ONLY
\newcommand{\AnalogAlpha}{alpha particle}     % APPENDIX X ONLY
\newcommand{\AnalogNuclear}{nuclear}          % APPENDIX X ONLY

% ============================================================
% END PREAMBLE
% ============================================================
